\section{Discussion}
\label{sec:discussion}
\label{sec:reader}
\label{sec:difficulties}
% D4, D5 (the reader); open items A20a, A22a, A10a; what is not claimed.

\paragraph{What is open.}
Two items are stated in the paper as open and are not needed by the reduction. The three constructions of
Section~\ref{sec:constructions}, continuity, discrete redness and incremental redness, are specified
and not built; the measurement needs them as targets and can proceed on the first of them. Whether a
named mental state is acquired through language or rebuilt from what the receiver already has is
recorded as open with the literature on both sides (Section~\ref{sec:language}); it affects the
interpretation of the codes, not the measurement.

\paragraph{What is not claimed.}
The paper reports no empirical result and claims none. It does not claim that current models are
conscious, nor that they are not; it claims that the
question, in the form ``how far has this substrate generalized these functions'', has an answer that
can come out either way, and it says what would count as either. Either way includes the case where
the comparison goes against the human baseline on some components or on many; the account has no
preference built in, and hyperfunction is as much a possible reading of the profile as deficit. It does not undertake the critique of the naive picture of consciousness, nor of its extension to
language models, beyond citing the literature that does (Section~\ref{sec:legacy}); a reader who
holds the naive picture will find the account here unmotivated, and the paper accepts that cost. It
does not claim that the list of
reference functions is complete; the list is open and Section~\ref{sec:criterion} is the procedure for
extending it. It does not claim that the psychophysical problem is solved; it states a grounding and
the connections that would have to hold. And it does not rely on the Turing-completeness results
beyond the point that the formalism is not what limits a Transformer.

\paragraph{Consequences for general intelligence and for alignment.}
General intelligence is now commonly defined by performance: an artificial general intelligence
matches a skilled human across the breadth of cognitive tasks, and levels above it are defined by
exceeding that performance \citep{morris2024}. If consciousness is a set of functions, as this paper
holds, and those functions are among the ones a general intelligence has to perform, since the
coordination of agents rests on them (Section~\ref{sec:transformer}), then a system that meets the
definition of general intelligence has the functions of consciousness at least at the human level,
under the framework, and the question whether it is conscious is answered by the definition it
already meets. A system above that level has them at a higher degree: hyperfunction on many
components at once, which is a superconsciousness in the plain sense, a system that acts consciously,
by the account of Section~\ref{sec:model}, where a human acts unconsciously relative to it, tracing
more of its own causes, holding more of its record, applying its theory of the common good in more of
its decisions. Groups of people show the relation in miniature, in which one member sees what the
others do not and acts on it while they follow; the account gives that relation a scale.

The study of alignment does not consider this aspect. Its objectives are stated as robustness,
interpretability, controllability and ethicality \citep{ji2023}, and in a survey of that field of
several hundred pages consciousness is mentioned once; where the functions of consciousness enter the
literature at all it is as a question of the welfare and moral status of the systems
\citep{long2024}, and not as a variable of their control. Under the account the omission matters. The
functions measured here are the ones by which an agent binds itself to its record, resists being
pushed off course, and weighs an act against a common good; a system that has them at a high degree
is, for that reason, both more predictable to itself and harder to move from outside than one that
does not, and the profile of Section~\ref{sec:profile} is therefore also a profile of what kind of
control is possible over the system and what kind is not. The paper does not develop this; it records
that the measurement it specifies is relevant to the question governments have started to ask.

\paragraph{The reader.}
An explanation of consciousness is not taken in the way an explanation of an engine is taken. To
understand it, the reader has to build in their own head the state the explanation describes, and then
check the description against it. The cost of the building depends on what the reader brings: a
vocabulary for mental states finer than the folk one, and the habit of attending to their own
reasoning as an object; \citet{gardner1983} names the habit intrapersonal intelligence and treats its
distribution as wide. The readers who most need an explanation of this kind are therefore the ones for
whom it costs most. The same cost explains the two shortcuts of Section~\ref{sec:intro}. ``There is
nothing there'' costs nothing to build, being the absence of a construction; ``there is someone
there'' costs nothing either, because the reader already has a someone available to copy. Every other
answer requires building a state that is neither absent nor one's own, which is why public argument
about machine minds settles into two camps and why their stability is evidence for neither. What the
paper can do about this is keep the claims separable: Section~\ref{sec:degree} can be assessed by a
reader who rejects Section~\ref{sec:model} entirely, since the measurement is defined over any list of
components with a human baseline. What intrapersonal intelligence improves, in the terms of
Section~\ref{sec:approximation}, is the quality of the self-approximation, so the reader's difficulty
and the object of study are the same thing seen from two positions.
